\documentclass[a4paper,12pt]{article}
\usepackage{amsmath}
\usepackage{graphicx}
\usepackage{geometry}
\usepackage{xcolor}
\geometry{top=1in,bottom=1in,left=1in,right=1in}

\title{ME473 Dynamic Finite Element Analysis Mini-Project}
\author{Stefano Burzio}
\date{}

\begin{document}

\maketitle

\section*{Evaluation Criteria}

The evaluation of the mini-projects will be based on the following key aspects. The criteria highlighted in blue apply exclusively to mini-project 3.

\subsection*{1. Problem Definition and Objectives (15\%)}
\begin{itemize}
    \item Clear identification and understanding of the problem statement.
    \item {\color{blue} Well-defined project objectives and scope}.
    \item {\color{blue} Relevance and importance of the chosen topic in the context of dynamic FEA.}
\end{itemize}

\subsection*{2. Methodology and Approach (15\%)}
\begin{itemize}
    %\item Appropriateness and correctness of the chosen FEA method (e.g., Modal, Transient, or Harmonic).
    \item Proper setup of material properties, geometry, and boundary conditions.
    \item Accuracy in the definition of the nodes and the meshing.
    \item Clear explanation of assumptions and simplifications made during the analysis.
\end{itemize}

\subsection*{6. Coding and Simulation (20\%)}
\begin{itemize}
    \item Code organization, readability, and commenting for understanding.
    \item Proper use of libraries or functions for solving dynamic FEA problems.
    \item {\color{blue} Successful integration of the code with the FEA software (if applicable), or successful standalone solution.}
\end{itemize}

\subsection*{3. Results and Analysis (20\%)}
\begin{itemize}
    \item Clear presentation and interpretation of results.
    \item Comparison of simulated results with theoretical values (if available) or expected outcomes.
    \item Identification and discussion of results (e.g., resonance, natural frequencies, peak responses).
    \item {\color{blue} Completeness of results (e.g., mode shapes, frequency analysis, time-history plots, displacement, stress distributions).}
\end{itemize}


\subsection*{4. Report Structure and Presentation (20\%)}
\begin{itemize}
    \item Clear, well-organized report structure with appropriate sections (Introduction, Methodology, Results, etc.).
    \item Quality of visuals (figures, graphs, screenshots and tables) and their relevance to the analysis.
    \item Proper citations and references. Plagiarism is not allowed and will be sanctionned. The use of AI tools is permitted, provided that you clearly cite the reference and specify how the tool was utilized in your work.
\end{itemize}

\subsection*{5. Innovation and Critical Thinking (10\%)}
\begin{itemize}
    \item Creativity in solving the problem or proposing alternative approaches.
    \item Depth of analysis and the ability to critically discuss the results.
    \item Suggestions for future work or improvements in the modeling and analysis process.
\end{itemize}



\end{document}
